\documentclass[11pt]{article}
\usepackage[T1]{fontenc}
\usepackage{lmodern}
\usepackage[margin=1in]{geometry}
\usepackage{microtype}
\usepackage{amsmath,amssymb,amsthm}
\usepackage{booktabs}
\usepackage{url}
\usepackage[hidelinks]{hyperref}
\newtheorem{theorem}{Theorem}
\newcommand{\pinn}{\operatorname{pinn}}
\newcommand{\leM}{\leq_{\mathrm M}}
\newcommand{\doi}[1]{\href{https://doi.org/#1}{\nolinkurl{doi:#1}}}
\hypersetup{
  pdftitle={A counterexample to distributivity of fixed pinnacle-order fibers},
  pdfauthor={Xiaojian Zeng}
}
\title{A counterexample to distributivity of\protect\\ fixed pinnacle-order fibers}
\author{Xiaojian Zeng\\[3pt]
  \small PBC School of Finance, Tsinghua University, 100190\\
  \small \href{mailto:xiaojianzeng@foxmail.com}{\texttt{xiaojianzeng@foxmail.com}}}
\date{}

\begin{document}
\maketitle

\begin{abstract}
Liang and Sagan asked whether permutations with a prescribed left-to-right
ordering of their pinnacle values form a distributive lattice in the
middle order. We give a negative answer using permutations of $[5]$ whose
only pinnacle is $5$. Five explicit permutations violate a distributive
identity, with every required meet and join determined in the entire
fiber. The decisive meet is certified by listing four possible inversion
tables. Thus the counterexample concerns a nonempty fiber and has a
self-contained proof independent of exhaustive computation.
\end{abstract}

\section{The question}

Write $S_n$ for the permutations of $[n]=\{1,\ldots,n\}$ in one-line
notation. An interior position $i$ is a peak of $\pi\in S_n$ when
$\pi_{i-1}<\pi_i>\pi_{i+1}$, and the value $\pi_i$ is a pinnacle
\cite{DNPT}. Endpoints are excluded. Let $\pinn(\pi)$ be the complete
sequence of pinnacle values, read from left to right, and define
\[
  P_n(\sigma)=\{\pi\in S_n:\pinn(\pi)=\sigma\}
\]
for a sequence $\sigma$ of distinct positive integers. This records the
actual pinnacle values in their specified order.

For $v\in[n]$, set
\[
 I_v(\pi)=\#\{u<v:\pi^{-1}(u)>\pi^{-1}(v)\},
 \qquad I(\pi)=(I_1(\pi),\ldots,I_n(\pi)).
\]
The middle order is
\[
 \pi\leM\rho\quad\Longleftrightarrow\quad
 I_v(\pi)\le I_v(\rho)\text{ for every }v\in[n].
\]
These are coordinates indexed by values, as in
Bouvel--Ferrari--Tenner \cite{BFT} and Liang--Sagan \cite[Section~6]{LS}.
The full middle order is a product of chains and hence a distributive
lattice \cite[Section~1]{BFT}.

For completeness, the coordinate bijection has an elementary description.
Given integers $0\le t_v\le v-1$, insert $v=1,2,\ldots,n$ successively
into the current word, placing $v$ in the unique slot with exactly $t_v$
entries to its right. Later insertions involve larger values and preserve
the count of smaller entries to the right of $v$. This constructs the
unique permutation with $I_v=t_v$. Consequently, ambient meets and joins
correspond to coordinatewise minima and maxima.

Liang and Sagan \cite[Question~8.6(a)]{LS} ask whether the induced middle
order on $P_n(\sigma)$ is a distributive lattice for every $n$ and
$\sigma$. The same question is numbered 8.7(a) in their first arXiv
version. Rusu and Tenner \cite{RT} characterize which pinnacle orderings
can occur, through a criterion limiting interruptions by larger
pinnacles. That existence criterion does not determine meets and joins
among the realizing permutations. The counterexample below supplies
such internal operations and shows that ambient distributivity is lost
on a fixed ordered fiber.

The question can also be restricted explicitly to nonempty fibers,
avoiding any convention about whether an empty poset is a lattice.
Our example answers this qualified question negatively as well.

\section{The counterexample}

\begin{theorem}\label{thm:main}
The nonempty fiber $P_5((5))$, with its induced middle order, is not a
distributive lattice.
\end{theorem}

\begin{proof}
Put $P=P_5((5))$. All meets and joins used below are bounds in the
whole of $P$. Consider the five permutations
\[
\begin{array}{@{}cccl@{}}
\toprule
\text{Name}&\text{Permutation}&I(\pi)&\text{Only peak position}\\
\midrule
e&12354&(0,0,0,0,1)&4\\
b&12534&(0,0,0,0,2)&3\\
a&13524&(0,0,1,0,2)&3\\
c&13452&(0,0,1,1,1)&4\\
d&13542&(0,0,1,1,2)&3\\
\bottomrule
\end{array}
\]
Each has precisely one pinnacle, of value $5$, so all five lie in $P$.
Their coordinates give $e<_{\mathrm M}b<_{\mathrm M}a<_{\mathrm M}d$
and $e<_{\mathrm M}c<_{\mathrm M}d$. In particular, the operations on
the comparable pairs satisfy
\begin{equation}\label{eq:comparable}
 a\wedge d=a,\qquad a\wedge b=b,\qquad b\vee e=b.
\end{equation}
The coordinate maximum of $I(b)$ and $I(c)$ is $I(d)$. Since $d\in P$,
it is an upper bound in $P$, and every upper bound in $P$ has each
coordinate at least the corresponding coordinate of $d$. Therefore
\begin{equation}\label{eq:join}
 b\vee c=d.
\end{equation}

It remains to find $a\wedge c$. Every ambient common lower bound has
inversion table
\[
 (0,0,u,0,v),\qquad u,v\in\{0,1\}.
\]
The coordinate bijection shows that the following four permutations
exhaust these bounds:
\[
\begin{array}{@{}ccc@{}}
\toprule
I(\pi)&\pi&\pinn(\pi)\\
\midrule
(0,0,0,0,0)&12345&()\\
(0,0,0,0,1)&12354&(5)\\
(0,0,1,0,0)&13245&(3)\\
(0,0,1,0,1)&13254&(3,5)\\
\bottomrule
\end{array}
\]
Only $e=12354$ belongs to $P$. Thus $e$ is the only common lower bound
of $a,c$ in $P$, and in particular
\begin{equation}\label{eq:meet}
 a\wedge c=e.
\end{equation}
This establishes the internal meet even though the ambient meet,
$13254$, is outside the fiber.

Equations \eqref{eq:comparable}--\eqref{eq:meet} now give
\begin{align*}
 a\wedge(b\vee c)&=a\wedge d=a,\\
 (a\wedge b)\vee(a\wedge c)&=b\vee e=b.
\end{align*}
Since $a=13524\ne12534=b$, the distributive identity fails. Every
operation in this identity has been justified within $P$; if $P$ were
a distributive lattice, these operations would have to satisfy the
identity. This contradiction proves the theorem.
\end{proof}

\section{Verification and scope}

An exact enumeration of the $120$ permutations of $[5]$ gives $56$
elements of $P$. The computation checked the common lower and upper
bounds for all $25$ ordered pairs of the five named elements against
the entire fiber. A separately written verifier reconstructed the
permutations by retaining the words with distinct entries in
$\{1,\ldots,5\}^5$, computed inversion coordinates from inverse
positions, and obtained the same fiber and bounds. These checks use
only integers and finite comparisons.

Zeng et al.\ \cite[Proposition~1 and Section~6]{ZLRO} likewise use exact
predicates and separately reconstructed finite geometric models to
separate computation from certificate checking. Here the two displayed
tables and the bound arguments already give a complete certificate.

To reproduce the finite check, list the permutations with pinnacle
sequence $(5)$ and compare their five inversion coordinates. For each
named pair $x,y$, form
\[
 L_{x,y}=\{z\in P:z\leM x,\ z\leM y\},\qquad
 U_{x,y}=\{z\in P:x\leM z,\ y\leM z\}.
\]
A proposed meet $m$ is verified by checking $m\in L_{x,y}$ and
$z\leM m$ for every $z\in L_{x,y}$; verify a proposed join dually.
No floating-point tolerance, random sampling, symmetry reduction, or
proof assistant is involved.

Theorem~\ref{thm:main} gives a negative answer to the universal assertion
in Question~8.6(a) of \cite{LS}. It makes no assertion of minimality in
$n$, no classification of other fibers, and no claim concerning the
separate counting question 8.6(b). A focused literature comparison on
14~September~2026 found no equivalent counterexample in the inspected
sources; this is a bounded assessment rather than a claim of priority.

\paragraph{Computational and AI assistance.}
OpenAI Codex assisted with literature searches, counterexample
exploration, proof development, exact verification code, a separate
mathematical review pass, and manuscript preparation. The separate
model context is not external peer review. The mathematical basis is
the explicit argument above, corroborated by the finite computations.

\bibliographystyle{abbrvurl}
\bibliography{references}
\end{document}
